% 原书第 91—94 页；第 6.6.4—6.6.5 节。
\subsection{缺失数据模拟}
\label{subsec:ch06-missing-data}

回顾一下，HAL 的一个突出特点是：无论维数 $d$ 为多少，都能保证其收敛速率快于 $n^{-1/4}$。这个速率恰好是保证所得 TMLE 有效性时，干扰参数初始估计所需的临界速率。因此，我们很希望了解，这个出色的渐近结果，究竟在多大程度上能使有限样本中的 TMLE 估计量表现良好。为研究这一问题，我们考虑了与预测模拟相同的九个数据生成分布，但在数据结构中进一步引入了缺失，令 $A=1$ 表示结局 $Y$ 被观测到。缺失按如下方式模拟：设 $g_0(w)=P(A=1\mid W=w)$，在“光滑型”情形下，
% 源书 g_0^5 的定界符为 {( ... ))，按印刷保留。
\[
 \begin{aligned}
 g_0^1(W)={}&\expit\{(w_1+4)^{1/2}-w_1/2\},\\
 g_0^3(W)={}&\expit\{1+(w_1+4)^{1/2}-w_1/2+w_2w_1/5-w_2^2/10\},\\
 g_0^5(w)={}&\expit\{(1+(w_1+4)^{1/2}-w_1/2+w_2w_1/5-w_2^2/10\\
           &\qquad{}+w_4+w_5/5-w_4w_3w_1/5));
 \end{aligned}
\]
在“跳跃型”情形下，
% 源书 g_0^1 以 )) 而非 )} 结束；各式系数、顺序、阈值按原书保留。
\[
 \begin{aligned}
 g_0^1(W)={}&\expit\{-3+2I(w_1<-3)-1.5I(w_1\geq-3,w_1<-1.5)\\
           &\qquad{}+0.5I(w_1\geq-1.5,w_1<0.5)-2I(w_1\geq0.5,w_1<2)\\
           &\qquad{}+2.4I(w_1\geq2)),\\[.5em]
 g_0^3(w)={}&\expit\{0.1I(w_1<-3)w_3+I(w_2<0)2.5\\
           &\qquad{}+1.5I(w_1\geq-3,w_1<-1.5)
                       +2I(w_1\geq-1.5,w_1<0.5)\\
           &\qquad{}-0.8I(w_1\geq0.5,w_1<2)+0.75I(w_1\geq2)w_3\\
           &\qquad{}+w_3-2I(w_1<0,w_2>0)\},\\[.5em]
 g_0^5(w)={}&\expit\{1+0.1I(w_1<-3)w_3+I(w_2<0)2.5\\
           &\qquad{}+1.5I(w_1\geq-3,w_1<-1.5)
                       +2I(w_1\geq-1.5,w_1<0.5)\\
           &\qquad{}-0.8I(w_1\geq0.5,w_1<2)+0.75I(w_1\geq2)w_3\\
           &\qquad{}+w_3-2I(w_1<0,w_2>0)+I(w_4<-1)\\
           &\qquad{}+2I(w_4<-2)-3I(w_4>0)+2I(w_5<3)\\
           &\qquad{}-1.7I(w_4<0)w_3\};
 \end{aligned}
\]
在“正弦型”情形下，
% 原书后两式在左花括号后多印一个未闭合的左圆括号，照录。
\[
 \begin{aligned}
 g_0^1(W)={}&\expit\{2+\sin(w_1)\},\\
 g_0^3(w)={}&\expit\{(2+\sin(w_1w_3)+\cos(w_2w_3)+\sin(w_1w_2)\},\\
 g_0^5(w)={}&\expit\{(1.5+\sin(w_1w_3)+\cos(w_2w_3)+\sin(w_1w_2)\\
           &\qquad{}+\sin(|w_4|)-w_3\cos(w_5^{1/2})\}.
 \end{aligned}
\]

对于九个数据生成分布中的每一个，我们在样本量为 500、2000 和 5000 时，各生成 500 次重复模拟，并使用基于不同干扰参数估计量的 TMLE，估计 $\psi_0=E_0\{E_0(Y\mid A=1,W)\}$。具体而言，我们考虑用预测模拟中相同的非参数估计量来估计 $\bar Q_0$ 和 $g_0$，也考虑以这些估计量为候选算法的超级学习器和离散超级学习器。我们关注的是，各种 TMLE 在有限样本中，其抽样分布（sampling distribution）在多大程度上近似正态。为此，我们根据 TMLE 的真实渐近方差，计算 Wald 型 95\% 置信区间的覆盖率，并通过视觉检查抽样分布的直方图来评价。

\begin{figure}[htbp]
 \centering
 % 原始图中各数据点、参照线、刻度和顺序全部保留；仅替换文字标签。
\resizebox{.88\textwidth}{!}{%
\begin{tikzpicture}[x=1bp,y=-1bp]
 \node[anchor=north west,inner sep=0pt] at (0,0) {\includegraphics[width=294bp,height=183bp]{figures/ch06_fig5_source.pdf}};
 \fill[white] (0,0) rectangle (88,159);
 \fill[white] (125,172) rectangle (260,183);
 \node[font=\fontsize{8.0}{9}\selectfont] at (192.1,177.9) {TMLE 的名义 95\% 覆盖率};
 \node[anchor=east,inner sep=0pt,font=\fontsize{8.0}{9}\selectfont] at (85.1,146.739) {径向 SVM};
 \node[anchor=east,inner sep=0pt,font=\fontsize{8.0}{9}\selectfont] at (85.1,125.799) {核回归};
 \node[anchor=east,inner sep=0pt,font=\fontsize{8.0}{9}\selectfont] at (85.1,105.731) {回归树};
 \node[anchor=east,inner sep=0pt,font=\fontsize{8.0}{9}\selectfont] at (85.1,86.550) {GBM};
 \node[anchor=east,inner sep=0pt,font=\fontsize{8.0}{9}\selectfont] at (85.1,65.729) {离散超级学习器};
 \node[anchor=east,inner sep=0pt,font=\fontsize{8.0}{9}\selectfont] at (85.1,46.360) {随机森林};
 \node[anchor=east,inner sep=0pt,font=\fontsize{8.0}{9}\selectfont] at (85.1,25.615) {超级学习器};
 \node[anchor=east,inner sep=0pt,font=\fontsize{8.0}{9}\selectfont\bfseries] at (85.1,6.286) {HAL};
\end{tikzpicture}%
}

 \caption{基于不同干扰参数估计量的 TMLE 估计量，其 Wald 型置信区间的覆盖率。结果按覆盖率与 95\% 之间的平均绝对距离排序。}
 \label{fig:ch06-tmle-coverage}
\end{figure}

图 6.5 展示了 27 种不同模拟情形中 Wald 型置信区间的覆盖率。使用 HAL 估计干扰参数的 TMLE 表现出色：其覆盖率平均只偏离名义 95\% 覆盖率约 1.1 个百分点，而且在任何模拟中都不低于 90\%。基于超级学习器的 TMLE 也产生了覆盖率很好的置信区间；不过，在一元“跳跃型”情形中，样本量较小时，基于 HAL 的 TMLE 明显更好。基于 GBM 的 TMLE 也有类似情况：除了两种“跳跃型”情形外，它在所有情形中的覆盖率都很好；而那两种情形的覆盖率很差，拖累了这类估计量的总体表现。

HAL 快速收敛速率的益处，在图 6.6 的直方图中清晰可见。该图展示了五元“跳跃型”情形中，样本量分别为 500 和 5000 时 TMLE 的抽样分布。第一行表明，即使样本量较小，基于 HAL 的 TMLE 也已近似正态，且偏差很小。样本量较大时，基于 HAL 的 TMLE 几乎没有偏差，其抽样分布可由图示正态分布很好地近似。相比之下，基于核回归的 TMLE 在小样本中存在严重偏差，而且我们可以清楚地看到，这一偏差并没有以快于 $n^{-1/2}$ 的速率收敛到零。

\begin{figure}[htbp]
 \centering
 % 图 6.6 的直方图为原书嵌入的 575×381 图像；保留原数据，替换核回归行标签。
\resizebox{.88\textwidth}{!}{%
\begin{tikzpicture}[x=1bp,y=-1bp]
 \node[anchor=north west,inner sep=0pt] at (0,0) {\includegraphics[width=276bp,height=182.52bp]{figures/ch06_fig6_source.pdf}};
 \fill[white] (0,94) rectangle (9.5,176);
 \node[rotate=90,font=\fontsize{8}{9}\selectfont] at (4.5,137) {核回归};
\end{tikzpicture}%
}

 \caption{基于 HAL 和基于核回归的标准化 TMLE 估计量的抽样分布。实线表示有效估计量的渐近分布；各估计量的均值标在相应的横轴上。}
 \label{fig:ch06-tmle-distribution}
\end{figure}
\FloatBarrier

\subsection{结论}
\label{subsec:ch06-conclusion}

本节考察了 HAL 估计量的实际构造和表现。我们发现，无论是用于预测，还是用于估计 TMLE 所需的相关中间干扰参数，它的表现都十分出色。这些结果表明，HAL 对构建稳健的超级学习器算法库具有重要价值，在实践中纳入它，很可能带来实际收益。关于 HAL 的早期结果，见会议论文 Benkeser and van der Laan（2016）。

% 清出本章末页后恢复外层的页底设置。
\clearpage
\endgroup
